% formal/random-polytope-boundedness.tex — input by formal/main.tex
%
% Boundedness characterization for rejection sampling of random 4-polytopes.
% Contains the sampling procedure. Proposition~\ref{prop:boundedness-iff}
% and its proof now live in exact-polar-geometry.tex.
%
% Sources:
%   - Proposition and proof: standard convex geometry (recession cones,
%     positive span). Original proof written by Claude, reviewed by Jörn;
%     see the corrected endpoint step and status in exact-polar-geometry.tex.
%   - Empirical acceptance rates: topics/empirical-studies/sys-landscape/rejection-calibration/
%
% Structure:
%   - Sampling procedure (halfspace representation, 4-step validation)
%   - Proposition: bounded iff (span R^4) and (all triples pass)
%   - Proof: recession cone ↔ bounded, positive span ↔ (i)+(ii)

% Jörn: original text approved (57e73d9). The vertex-enumeration sentence
% below was corrected for nonsimple vertices here; the boundedness proof in
% exact-polar-geometry.tex has a corrected interval-endpoint step. Both
% corrections await Jörn's review.
\subsection{Rejection sampling acceptance rates}
\label{sec:rejection-calibration}

To probe Viterbo's conjecture computationally, we need large datasets of
random convex polytopes in~$\mathbb{R}^4$. We generate these by rejection
sampling: sample a candidate halfspace representation, then validate that it
defines a bounded, irredundant polytope.

\subsubsection*{Sampling procedure}

A polytope~$K \subset \mathbb{R}^4$ is specified by dual vertices
$a_1, \ldots, a_F \in \mathbb{R}^4$:
\[
  K = \bigl\{\, x \in \mathbb{R}^4 \;\big|\; \langle a_i, x \rangle \le 1,\;
  i = 1, \ldots, F \,\bigr\}.
\]
Each dual vertex decomposes as $a_i = n_i / h_i$ where~$n_i \in S^3$ is the
outward unit normal and~$h_i > 0$ is the height
(positive heights ensure the origin lies in the interior).

For a given facet count~$F$ and height range~$[h_{\min}, h_{\max}]$, each
sample attempt proceeds as follows:
\begin{enumerate}
  \item Sample~$F$ independent unit normals uniformly on~$S^3$
    (via normalized standard Gaussian vectors in~$\mathbb{R}^4$).
  \item Sample~$F$ independent heights uniformly from~$[h_{\min}, h_{\max}]$.
    Set $a_i = n_i / h_i$.
  \item Validate the candidate~$(a_1, \ldots, a_F)$:
  \begin{enumerate}
    \item \emph{No duplicate normals}: reject if~$\|n_i - n_j\| < \varepsilon$
      for any~$i \ne j$.
      % Same-direction only (n_i ≈ n_j). Antiparallel normals (n_i ≈ -n_j)
      % are valid opposite-facing facets, e.g. hypercube has ±e_i.
    \item \emph{Boundedness}: the polytope must be bounded
      (Proposition~\ref{prop:boundedness-iff}). We first check that the
      dual vertices span~$\mathbb{R}^4$ (via SVD rank), then enumerate all
      triples~$(i,j,k)$: the three dual vertices share a 1D kernel direction~$d$
      (computed via the generalized cross product in~$\mathbb{R}^4$), and we
      verify that some $a_\ell$ has positive inner product with~$d$ and some
      has negative inner product with~$d$.
      % The degenerate case (three coplanar dual vertices, kernel ≥ 2D) is skipped in
      % the code — probability zero for continuous distributions on S^3.
    \item \emph{Vertex enumeration}: each vertex has at least~$4$ linearly independent active facet
      constraint hyperplanes, including nonsimple vertices. Solve all~$\binom{F}{4}$ linear systems,
      discard infeasible or duplicate solutions.
    \item \emph{Irredundancy}: each facet must be a genuine 3-dimensional
      face. Check that the vertices lying on each facet hyperplane span a
      3-dimensional affine subspace (via SVD).
  \end{enumerate}
  \item If all checks pass, accept the polytope.
\end{enumerate}
